\documentclass[11pt,a4paper]{article}
\usepackage{arxiv}

\title{Securing Neural Interfaces: Architecture, Threat Taxonomy, and Neural Impact Scoring for Brain-Computer Interfaces}

\author{
  Kevin L. Qi\\
  Qinnovate\\
  \texttt{qinnovate.com}
}

% Registry figures, generated from datalake/qtara-registrar.json.
% Snapshot* = the frozen v1.5 state this paper's analysis was computed against.
% Live*     = the registry as it stands at build time.
% Generated by research/paper/scripts/generate_registry_stats.py
% Do not edit. Run `make stats` in research/paper/ to regenerate.
%
% Snapshot* are the frozen registry v1.5 figures the paper's analysis was
% computed against. Live* are read from datalake/qtara-registrar.json at
% build time. Do not use a Live* macro as the denominator of a numerator
% that was computed over the snapshot.

\newcommand{\SnapshotRegistryVersion}{1.5}
\newcommand{\SnapshotTechniques}{109}
\newcommand{\SnapshotTactics}{15}
\newcommand{\SnapshotDomains}{8}
\newcommand{\SnapshotGapTechniques}{108}
\newcommand{\SnapshotGapPercent}{99.1}
\newcommand{\LiveRegistryVersion}{4.3}
\newcommand{\LiveTechniques}{173}
\newcommand{\LiveVariants}{1}
\newcommand{\LiveTactics}{17}
\newcommand{\LiveWithCvss}{103}
\newcommand{\LiveNissScored}{129}
\newcommand{\LiveNissProvisional}{10}
\newcommand{\LiveNissUnscored}{34}
\newcommand{\LiveGapTechniques}{103}
\newcommand{\LiveGapPercent}{59.5}
\newcommand{\LivePinsFlagged}{37}


\date{October 2026 (v1.6; first posted February 2026)}

\begin{document}

\maketitle

\begin{abstract}
Brain-computer interfaces (BCIs) are moving from research tools toward commercially
deployed medical devices, with companies including Neuralink, Synchron, Blackrock
Neurotech, and Paradromics advancing toward regulatory approval and new entrants such as
Merge Labs raising \$252M in seed funding~\cite{mergelabs2026}. These devices read and,
increasingly, write neural signals. We are unaware of a published security framework that
scores the harms specific to such devices: the Common Vulnerability Scoring System
(CVSS v4.0) has no metric for tissue damage, for unauthorised reading or disruption of
neural function, for violated consent, or for effects that cannot be reversed.

We present a proposed, unvalidated framework with four parts:
(1)~an \textbf{addressing model} that tags what a device touches with eleven unordered
codes for neural structures, the electrode-tissue interface and device systems;
(2)~\textbf{TARA}, a catalogue of BCI attack techniques, each with a tactic, an evidence
status and a therapeutic analog where one exists;
(3)~\textbf{NISS}, an impact-only CVSS v4.0 extension with six neural-specific metrics; and
(4)~the \textbf{Neural Impact Chain}, which associates techniques with DSM-5-TR and
ICD-10-CM diagnostic category references for threat modeling, not as diagnostic claims.

The analysis is of a frozen snapshot of \SnapshotTechniques{} techniques (registry
v\SnapshotRegistryVersion{}). In that snapshot CVSS cannot express the primary impact of
59 techniques (54.1\%) and misses part of it for 30 more. Fifty-two techniques are in the
direct-association class across 29 diagnostic category codes, and 78 (71.6\%) have a recorded
therapeutic analog: many attack mechanisms are shared with therapies, and what separates the
two is consent, dosage and oversight. This revision (v1.6) replaces the earlier hourglass
stack with addresses, adds an evidence rule and three scoring states, states NISS's limits,
corrects case studies that did not match the registry, and argues why neurosecurity and
neuroethics are needed now. Live registry figures (v\LiveRegistryVersion{},
\LiveTechniques{} techniques) are reported beside the snapshot, not substituted into it.
The framework, catalogue and scoring data are released under the Apache 2.0 license.
\end{abstract}

\noindent\textbf{Keywords:} brain-computer interface, neurosecurity, CVSS, threat taxonomy,
neural scoring, neuroethics, DSM-5-TR, dual-use

\tableofcontents
\newpage

% ═══════════════════════════════════════════════════════════════
% Section 1: Introduction
% ═══════════════════════════════════════════════════════════════

\section{Introduction}
\label{sec:introduction}

Brain-computer interfaces (BCIs) are no longer confined to research laboratories.
Neuralink has implanted its N1 device in human patients~\cite{musk2019neuralink},
Synchron's Stentrode has demonstrated motor neuroprosthesis via neurointerventional
surgery~\cite{oxley2021stentrode}, and Blackrock Neurotech's Utah array has enabled
high-performance speech neuroprostheses~\cite{willett2023speech}. Paradromics is
advancing toward high-bandwidth cortical interfaces for paralyzed patients.
Investment in neurotechnology grew 700\% between 2014 and 2021, with the global
market projected to reach \$25 billion by 2030~\cite{unesco2025recommendation}.

Regulation has begun to require cybersecurity for these devices without saying how to
assess it. Section 3305 of the Food and Drug Omnibus Reform Act of 2022
(FDORA, Pub.~L. 117-328) added Section 524B to the FD\&C Act~\cite{fda524b},
mandating cybersecurity documentation---including threat modeling, software bills
of materials, and vulnerability disclosure---in premarket submissions for
connected medical devices, and since October 2023 the FDA has applied a
Refuse-to-Accept policy to non-compliant submissions. Section 524B does not specify
a threat taxonomy or scoring system, and we are unaware of provisions in the standards
it draws on (CVSS, IEC 62443, AAMI TIR57, ISO 14971) for the risks specific to devices
that read and write neural signals. We found no neural-specific security requirements in
the EU Medical Device Regulation~\cite{eumdr2017}.

\subsection{Why Neurosecurity}
\label{sec:why-neurosecurity}

Four properties of neural devices, each documented in the literature, separate their
security from that of other connected medical devices.

\paragraph{The device acts on a person.}
A stimulating device delivers current to tissue. The same mechanisms underlie established
therapies such as deep brain stimulation~\cite{lozano2019dbs} and transcranial magnetic
stimulation~\cite{hallett2007tms}, so an unauthorised change to stimulation parameters is a
change to a treatment, not to a record. Some such effects may not be reversible, and a
security assessment has to say so.

\paragraph{Neural data cannot be reissued.}
A leaked password can be changed; a recording of a person's brain activity cannot. Present
devices do not read thoughts~\cite{ienca2018brainleaks,wexler2019neurohype}: speech
neuroprostheses decode attempted speech using models trained on a cooperating
user~\cite{willett2023speech}. Information can nevertheless leak. Martinovic et
al.~\cite{martinovic2012feasibility} showed that P300 responses recorded with a consumer
headset reveal information about PINs, banks and other private details at better than
chance. If decoders keep improving, a recording retained today becomes a readout risk later; the
catalogue records this as harvest-now-decode-later (QIF-T0166), with theoretical status.

\paragraph{The device can outlast its support.}
An implant stays in a person after the party that maintains it withdraws. In a modified
Delphi study with a 24-expert panel, the 14 experts who completed the final survey agreed
that responsibility for implanted neural devices after a trial needs explicit
rules~\cite{higgins2026posttrial}. No attacker is needed for harm
here, which is why this revision keeps such conditions in a separate hazards dataset
(Section~\ref{sec:tara}).

\paragraph{The attack surface spans biology and silicon.}
A recent review treats every element of the loop between brain, device, linked devices and
servers, and every link between them, as an attack surface, and describes defence frameworks
built specifically for BCIs as research prototypes~\cite{bagley2026threatvectors}.

Two consequences follow. First, an assessment needs dimensions that information-technology
scoring lacks (Section~\ref{sec:niss}). Second, the evidence base is thin: of the
\LiveTechniques{} techniques in the live catalogue, 95 are theoretical, and none of the
catalogue's status values means that an attack has been carried out on a person with a
device. Neurosecurity at this stage is therefore anticipatory work, and its credibility
depends on separating what has been shown from what is supposed. The evidence rules added
in this revision (Section~\ref{sec:tara}) exist for that reason.

\subsection{The Scoring Gap}

The Common Vulnerability Scoring System (CVSS v4.0)~\cite{first2023cvss4} is the
industry standard for assessing vulnerability severity. CVSS captures exploitability
characteristics---attack vector, complexity, privileges required, user
interaction---and impact on the confidentiality, integrity and availability of
information systems. It has no metric for:

\begin{itemize}
  \item \textbf{Biological tissue damage}---seizure induction, tissue injury,
        involuntary motor activation
  \item \textbf{Unauthorised reading or disruption of neural function}---inference
        from neural data, perceptual or motor disruption
  \item \textbf{Consent boundaries}---the distinction between operating within
        consented parameters and covert manipulation
  \item \textbf{Damage reversibility}---IT assets can be restored from backup;
        neural tissue cannot
  \item \textbf{Neuroplastic consequences}---prolonged stimulation can cause
        lasting change to neural pathways
\end{itemize}

When a vulnerability in a BCI could induce seizures, expose information inferred from
neural signals, or cause lasting harm, a CVSS base score alone is insufficient. In the
\SnapshotTechniques{}-technique snapshot analysed here, CVSS cannot express the primary
impact of 59 techniques and misses part of it for 30 more (Section~\ref{sec:niss}).

\subsection{Contributions}

This paper presents a proposed, unvalidated framework with four parts:

\begin{enumerate}
  \item \textbf{An addressing model} (Section~\ref{sec:hourglass}): eleven codes in
        use---seven for neural structures (N1--N7), one for the electrode-tissue
        interface (I0) and three for device systems (S1--S3)---each with a membership
        rule. The codes are addresses with no order. Earlier versions of this paper
        presented them as an hourglass stack with a chokepoint; Section~\ref{sec:hourglass}
        says why that was withdrawn.

  \item \textbf{TARA} (Section~\ref{sec:tara}): a catalogue of
        \SnapshotTechniques{}~techniques across \SnapshotTactics{}~tactics in the analysed
        snapshot (\LiveTechniques{} in the live registry), each classified by evidence
        status, severity and dual-use therapeutic potential. TARA is an independent
        taxonomy inspired by the structure of MITRE ATT\&CK\textsuperscript{\textregistered}.

  \item \textbf{NISS} (Section~\ref{sec:niss}): the Neural Impact Scoring
        System, an impact-only CVSS v4.0 extension designed to conform with
        FIRST.org's extension mechanism (User Guide
        \S3.11)~\cite{first2023cvss4userguide}, adding six neural-specific
        metrics, with stated limits and three scoring states.

  \item \textbf{The Neural Impact Chain} (Section~\ref{sec:nic}): a threat-modeling
        method that associates techniques with diagnostic category references through
        addresses, structures and functions to DSM-5-TR and ICD-10-CM
        codes~\cite{apa2022dsm5tr}. The codes are category references, not
        diagnostic claims.
\end{enumerate}

Section~\ref{sec:governance} sets out why a security framework needs neuroethics and
how it enters the scoring; Section~\ref{sec:case-studies} presents case studies, and
Section~\ref{sec:limitations} states limitations and what this revision changed. The
framework, catalogue and scoring system are released under the Apache 2.0 license.

% ═══════════════════════════════════════════════════════════════
% Section 2: Related Work
% ═══════════════════════════════════════════════════════════════

\section{Related Work}
\label{sec:related-work}

Research at the intersection of cybersecurity and neurotechnology has progressed
through three phases: foundational framing, empirical demonstration, and emerging
policy response. \qif builds on each while addressing gaps left by prior work.

\subsection{Foundational Neurosecurity}

Denning, Matsuoka, and Kohno~\cite{kohno2009neurosecurity}
introduced the term ``neurosecurity'' and analyzed attack
surfaces in implantable neurostimulators. Their work identified wireless
reprogramming, battery depletion, and signal injection as threat categories, but
predated modern high-bandwidth BCIs and did not propose a scoring system.

Bonaci, Calo, and Chizeck~\cite{bonaci2015appstores} extended this with ``App
Stores for the Brain,'' examining privacy threats from third-party BCI
applications. They identified the need for access control at the neural data
layer---a concern our governance framework addresses through consent tiers
(Section~\ref{sec:governance}).

\subsection{Empirical Attack Demonstrations}

Martinovic et al.~\cite{martinovic2012feasibility} showed at USENIX Security
2012 that P300 event-related potentials recorded with consumer-grade EEG headsets
leak information about PINs, banks and other private details at better than chance.
TARA (Section~\ref{sec:tara}) records this as P300 interrogation (QIF-T0035), with
status demonstrated.

Ienca and Haselager~\cite{ienca2016hacking} broadened the scope to ``hacking the
brain,'' analyzing BCI security through the lens of informational and physical
integrity. They proposed that BCI security requires both traditional cybersecurity
measures and novel neuroethical safeguards---a dual requirement that \qif
implements through the parallel CVSS/NISS scoring architecture.

Meng et al.~\cite{meng2023adversarial} established an adversarial robustness
benchmark for EEG-based BCIs, evaluating multiple adversarial defense approaches
across neural network architectures and EEG datasets. Their work focuses on
machine learning adversarial robustness; \tara extends coverage to physical-layer,
protocol-layer, and cognitive-layer threats.

\subsection{Emerging Frameworks}

Schroder et al.~\cite{schroder2025cyberrisks} (2025) developed a threat model for network-based attacks on next-generation BCIs
and identifying associated regulatory changes needed to address them. Their work
contributes threat categorization across network pathways but does not propose a
scoring system, formal threat taxonomy with technique-level granularity, or clinical
impact mapping.

Camara et al.~\cite{camara2015security} and Rushanan et al.~\cite{rushanan2014sok}
surveyed security and privacy in implantable medical devices more broadly, covering
pacemakers, insulin pumps, and cochlear implants alongside neurostimulators.
Halperin et al.~\cite{halperin2008pacemakers} demonstrated practical wireless
attacks against pacemakers, establishing that implanted medical device security is
not theoretical. These works address the broader medical device landscape; \qif
focuses specifically on the unique challenges of bidirectional neural interfaces.

\subsection{Neuroethics and Policy}

Ienca and Andorno~\cite{ienca2017neurorights} proposed four fundamental neurorights:
cognitive liberty, mental privacy, mental integrity, and psychological continuity.
Yuste et al.~\cite{yuste2017four} added equal access and protection from
algorithmic bias. These philosophical frameworks establish \emph{what} must be
protected; \qif provides the technical architecture for \emph{how}.

Muñoz et al.~\cite{lazaro2020researcher} and Lázaro-Muñoz et al.~\cite{lazaro2022posttrial} contributed
empirical grounding through researcher interviews, documenting clinician concerns
about adaptive DBS ethics (72\% cited uncertainty about risks) and post-trial
access obligations for implanted neural devices.

UNESCO's 2025 Recommendation on the Ethics of
Neurotechnology~\cite{unesco2025recommendation}---adopted by 194 Member
States---is a global normative instrument for neurotechnology governance. The OECD had previously published a Recommendation on Responsible
Innovation in Neurotechnology~\cite{oecd2019neurotechnology} covering 36 member
countries. Chile amended its constitution in 2021 to require the law to safeguard brain
activity and the information derived from it~\cite{chile2021neurorights}. These policy instruments establish
principles and values; \qif demonstrates technical implementation
(Section~\ref{sec:governance}).

\subsection{Cognitive Neurosecurity and Recent Surveys}

Two recent preprints from one group approach the problem from the cognitive side.
Bagley and Petritsch~\cite{bagley2024mathematicalframework} attribute the gap between
calls for neurotechnology security and applied work to the lack of a precise way to
describe the problems, and propose \emph{Cognitive Neurosecurity}: cognitive state
variables combined into a high-dimensional vector, with two attack classes, readout
attacks that identify cognitive states against a person's will and alteration attacks
that modify them. They state that the formalism borrowed from quantum theory implies no
physically quantum basis for cognition, treat defence as optimisation, expect defence
against alteration to depend on detection, and argue that open publication of attack
structures is the most ethical option. Their readout/alteration split corresponds to the
NISS split between Cognitive Reconnaissance and Cognitive Disruption
(Section~\ref{sec:niss}). They also note that small alterations spread across many people
could stay below individual-level detection; TARA records techniques one device and one
person at a time, so this is a gap in its model.

Bagley et al.~\cite{bagley2026threatvectors} review threat vectors and defences across
the loop from brain to servers. Mapping each vector in that review against TARA found
four without a corresponding technique; they were added as QIF-T0166, T0170, T0171 and
T0177, and one finding duplicated the published QIF-T0174 and is recorded there. Four of the five rest only on preprint
evidence and are unscored; QIF-T0166 has one indirect peer-reviewed source and carries a
provisional range. Both works are preprints, the mappings were
made with AI assistance, and their authors have not reviewed them.

\subsection{How QIF Differs}

Prior work has addressed BCI security, neuroethics and vulnerability scoring largely
separately. We are unaware of prior work that combines the following four elements in one
open dataset:

\begin{enumerate}
  \item \textbf{Addressing}: codes with membership rules for what a device touches
  \item \textbf{Taxonomy}: technique-level entries with an evidence status and an
        evidence rule (vs.\ category-level surveys)
  \item \textbf{Scoring}: a CVSS-compatible, impact-only extension with neural-specific
        metrics and explicit scoring states
  \item \textbf{Clinical association}: links from techniques to diagnostic category
        references for threat modeling
\end{enumerate}

We found no prior work that systematically associates vulnerability techniques with
DSM-5-TR diagnostic category codes, but our search was not exhaustive, and the
association method (Section~\ref{sec:nic}) has not been clinically validated.

% ═══════════════════════════════════════════════════════════════
% Section 3: From the Hourglass to an Addressing Model
% (label kept as sec:hourglass so earlier cross-references resolve)
% ═══════════════════════════════════════════════════════════════

\section{From the Hourglass to an Addressing Model}
\label{sec:hourglass}

Versions up to v1.5 of this paper presented an 11-band \emph{hourglass}: seven neural
bands (N7 neocortex down to N1 spinal cord), one interface band (I0, the electrode-tissue
boundary) and three synthetic bands (S1--S3), read as a stack. The model borrowed from the
OSI reference model~\cite{zimmermann1980osi} and from the Internet's narrow
waist~\cite{deering1998hourglass}, and made three claims: that I0 is a natural security
chokepoint through which all signals pass; that every component of a BCI maps to exactly
one band, so that no attack surface falls outside the model; and that behaviour becomes
less deterministic ascending the neural bands. This revision withdraws all three and
replaces the stack with addresses. This section records what changed and why.

\subsection{Why the Stack Was Withdrawn}

\paragraph{The numbering carries no flow.}
If the bands were a stack, neural pathways should mostly join numerically adjacent bands.
Recomputing each pathway's origin and target bands from the structure table gives 65
links between distinct N codes, of which 18 (27.7\%) join adjacent codes. Over all 5,040
possible numberings of the seven codes the mean is 28.6\%, and 60\% of numberings do at
least as well as the published one. The order is no better than chance.

\paragraph{The thalamus is not a gate.}
The v1.5 text used thalamocortical relay~\cite{sherman2006thalamus} to support
signal-plausibility checks, and earlier project documents described the thalamus as the
brain's firewall. The thalamus is the principal route by which sensory input, cerebellar
output and basal ganglia output reach neocortex, and transmission through it is
state-dependent~\cite{sherman2016thalamuscentral}. It is not a security control:
brainstem neuromodulatory systems also reach cortex by extra-thalamic
routes~\cite{edlow2012ascendingarousal}, motor thalamus integrates cortical, cerebellar and
basal ganglia input~\cite{boschbouju2013motorthalamus}, and olfaction is not thalamus-free,
since piriform cortex projects to the mediodorsal thalamus~\cite{courtiol2015olfactorythalamus}.

\paragraph{I0 is a boundary, not a chokepoint.}
The electrode-tissue boundary is where a recording or stimulating device transduces between
ionic and electronic current, and it remains an address. It is not a point through which
all attacks pass: of the \LiveTechniques{} techniques in the live catalogue, 110 carry no
I0 tag and 53 are tagged only at device codes, among them consumer-sensor side channels
that never touch an electrode.

\paragraph{The model was not complete.}
Nerves, roots, ganglia and sense organs had no address, so the completeness claim did not
hold. A periphery code (PER) is proposed for them and is not yet integrated.

\paragraph{Determinacy was untested.}
The per-band determinacy classes and ``QI'' ranges had no defined estimator. They are now
recorded as untested hypotheses (\texttt{datalake/qif-hypotheses.json}), proposed for the
earlier band memberships.

\subsection{Addressing Version 2}

The seven N codes are retained as identifiers. Each structure has one address, set by the
adult territory that holds its principal neurons; circuits such as the basal ganglia and
the limbic system are functional tags a structure can carry, not addresses.
Table~\ref{tab:bands} gives the codes in use. \textbf{The number carries no severity,
processing or traversal meaning.}

\begin{table}[H]
\centering
\caption{Address codes in use (addressing version 2). Listed by code, not by rank.}
\label{tab:bands}
\small
\begin{tabularx}{\textwidth}{@{}l l X@{}}
\toprule
\textbf{Code} & \textbf{Name} & \textbf{Membership rule} \\
\midrule
N7 & Neocortex       & Isocortex, including agranular motor areas \\
N6 & Allocortex      & Allocortex, mesocortex and pallial nuclei, including basolateral amygdala (v1.5: ``Limbic System'') \\
N5 & Cerebral nuclei & Subpallial nuclei: striatum, pallidum, basal forebrain, central and medial amygdala (v1.5: ``Basal Ganglia'') \\
N4 & Diencephalon    & Thalamus, subthalamus, hypothalamus, epithalamus \\
N3 & Cerebellum      & Cerebellar cortex and nuclei \\
N2 & Brainstem       & Midbrain, pons and medulla, including substantia nigra and ventral tegmental area \\
N1 & Spinal cord     & Spinal cord parenchyma; roots and ganglia are outside it \\
\midrule
I0 & Neural interface & Electrode-tissue boundary where biological signals become electronic \\
\midrule
S1 & Analog front-end   & Amplification, filtering, ADC/DAC conversion \\
S2 & Digital processing & Decoding algorithms, classification, digital signal processing \\
S3 & Application        & Clinical software, user interface, data storage, network \\
\bottomrule
\end{tabularx}
\end{table}

Three structures changed address. The substantia nigra and ventral tegmental area moved
from N5 to N2, because they lie in the midbrain, and the subthalamic nucleus moved from N5
to N4. All three keep their basal ganglia or reward tags. Subthalamic deep brain
stimulation, which v1.5 placed at N5, is therefore an N4 example. Because codes keep their
identifiers while membership changes, the atlas carries \texttt{addressing\_version: 2},
which should be cited with any band reference. Pathway bands are now computed from the
structure table by a script with a check mode, which the build does not yet run.

Fourteen statements in project documents drew security conclusions from the gate analogy,
chiefly that a deeper implant bypasses fewer protective layers. They are marked under
review, neither deleted nor asserted. The NISS formula has no term for implant depth: all
\LiveNissScored{} stored scores reproduce exactly from their six metric values, although
that does not show whether depth reasoning influenced the choice of any metric level.

\subsection{Routes (Proposed)}

Signal flow, which the stack could not carry, is moved to a proposed route view, staged and
not integrated. A route is an ordered list of synaptic hops, each with its sign,
transmitter and own evidence. A source check of all 90 hops in the 17 staged routes found
47 supported by a source a checker read, 42 partly supported and 1 contradicted for
primates; 3 hops are recorded as demonstrated in humans and 85 as inferred by homology from
other species. The routes are suitable for tagging what a device touches and for teaching.
They are not a basis for claims about what an attack can reach in a person.

\paragraph{Effect on this paper.}
The band annotations in the snapshot analysis (Sections~\ref{sec:tara}--\ref{sec:nic}) use
the version 1 memberships. Those analyses are not recomputed here; a reader using a band
count from them should treat it as a version 1 figure.

% ═══════════════════════════════════════════════════════════════
% Section 4: TARA Threat Taxonomy
% ═══════════════════════════════════════════════════════════════

\section{TARA: Threat Taxonomy}
\label{sec:tara}

The Therapeutic Applications \& Risk Assessment (TARA) registry catalogues BCI
attack techniques with a structure inspired by MITRE
ATT\&CK~\cite{mitre2024attack}, independently developed to cover neural,
cognitive, and physical safety domains. Each
technique is simultaneously an attack vector, an ethical risk, and---where
applicable---a therapeutic application.

\paragraph{Why not ATT\&CK directly?}
MITRE ATT\&CK is designed for enterprise IT and mobile environments where
adversary objectives center on data exfiltration, lateral movement, and
persistence. BCI threats differ in three fundamental ways that make direct
adoption inadequate: (1)~the target is biological tissue, not an information
system---attacks can cause seizures, tissue necrosis, or cognitive coercion,
none of which map to ATT\&CK's impact categories; (2)~the same technique is
built on the same mechanism as a therapy (e.g., the stimulation used in deep
brain stimulation for Parkinson's disease can also be delivered without authorisation), requiring a
dual-use classification that ATT\&CK has no mechanism to express; and (3)~the
attack lifecycle must span from radio-frequency transmission through silicon
processing to neural tissue, crossing the bio-digital boundary that ATT\&CK's
purely digital model does not address. TARA adopts ATT\&CK's proven structural
methodology---techniques organized into tactics and domains---while building an
independent taxonomy purpose-built for the BCI threat landscape.

\subsection{Taxonomy Structure}

The TARA registry organizes \SnapshotTechniques{} techniques into \SnapshotTactics{} tactics across \SnapshotDomains{} operational
domains. Technique identifiers follow the format \texttt{QIF-TXXX}, where the
numeric suffix provides sequential ordering.

\paragraph{Domains.}
The eight domains span the full attack lifecycle:

\begin{enumerate}
  \item \textbf{Neural (N)} --- Direct interaction with neural tissue
        (scan, injection, manipulation)
  \item \textbf{BCI System (B)} --- System-level intrusion and evasion
  \item \textbf{Protocol (P)} --- Protocol disruption and communication attacks
  \item \textbf{Data (D)} --- Data harvesting and exfiltration
  \item \textbf{Cognitive (C)} --- Cognitive exploitation and imprinting
  \item \textbf{Countermeasure (M)} --- Surveillance and monitoring
  \item \textbf{Evasion (E)} --- Defense evasion and anti-detection
  \item \textbf{Sensor (S)} --- Consumer device side-channel attacks
\end{enumerate}

\paragraph{Tactics.}
The 15 tactics follow a lifecycle structure analogous to ATT\&CK, adapted for BCI operations:
Neural Scan, BCI Intrusion, Neural Injection, Cognitive Imprinting, BCI Evasion,
Data Harvesting, Neural Manipulation, Evasion/Rootkit Deployment, Monitoring/Surveillance,
Protocol Disruption, Cognitive Exploitation, Signal Replay, Signal Harvesting,
Signal Fingerprinting, and Signal Chaining.

\subsection{Evidence Classification}

Each technique carries an evidence status. Table~\ref{tab:status} gives the snapshot
counts used in this paper beside the live counts.

\begin{table}[H]
\centering
\caption{TARA evidence status: the v\SnapshotRegistryVersion{} snapshot analysed here and
the live registry (v\LiveRegistryVersion{}, read 2026-10-17).}
\label{tab:status}
\small
\begin{tabularx}{\textwidth}{@{}l X r r@{}}
\toprule
\textbf{Status} & \textbf{Definition (live registry)} & \textbf{v\SnapshotRegistryVersion{}} & \textbf{Live} \\
\midrule
Confirmed    & Retired 2026-10-16 (see text)                                          & 19 & --- \\
Demonstrated & A published source executes the technique against a real system      & 35 & 48 \\
Emerging     & Shown against an analogous system or in part                         & 22 & 27 \\
Theoretical  & The mechanism is established but no attack using it has been shown   & 31 & 95 \\
Plausible    & The mechanism is contested or indirect; the attack is reasoned        & 1  & 2 \\
Speculative  & Neither mechanism nor attack is established                          & 1  & 1 \\
\midrule
             & \textbf{Total}                                                       & \SnapshotTechniques{} & \LiveTechniques{} \\
\bottomrule
\end{tabularx}
\end{table}

\paragraph{CONFIRMED retired.} The status \emph{confirmed} was retired on 2026-10-16
after an audit found that it spanned unrelated kinds of evidence; for example, it had been
given to signal injection (QIF-T0001) on the strength of a position paper and a P300
probing study, neither of which demonstrates signal injection. The 27 entries that carried
it in the live registry were reclassified as theoretical (21), demonstrated (5) and
plausible (1), each keeping its prior value and a reason. The 19 \emph{confirmed} entries of
the snapshot should be read with this in mind. \textbf{No status value means that an attack
has been carried out on a person with an implanted or worn device.}

\paragraph{Evidence rule.} Every live entry records whether its sources are peer-reviewed,
preprint or absent (158, 7 and 8 entries). \textbf{A technique with no peer-reviewed source
may not hold a status above theoretical and does not carry a NISS score.} Papers that test a
defence are held in a separate field, so they cannot satisfy the citation requirement for
an attack. When the rule was applied, 34 entries broke it; 24 were given a verified
peer-reviewed source and 10 lost their scores. ``Peer-reviewed'' currently means that no
preprint identifier was found, which is weaker than verified.

\paragraph{Hazards held separately.} Some conditions harm a device user with no attacker,
such as a manufacturer withdrawing support for an implant. They do not belong in a catalogue
of adversary actions, so 23 entries are held in \texttt{datalake/qif-hazards.json}. One,
QIF-H0001 (loss of manufacturer or sponsor support), rests on two peer-reviewed
sources~\cite{higgins2026posttrial,zhao2026consentinflux} and is theoretical. The other 22 are unsourced pointers
to catalogue records that describe a therapy or physical mechanism; they are not warnings
about those therapies and assert nothing about incidence in people.

\subsection{Severity Distribution}

Each technique also carries a filed severity rating. In v1.5 this table was labelled as
CVSS v4.0 base ratings; it is in fact the registry's \texttt{severity} field, assigned by
hand and not computed from the CVSS vectors. In the live registry it agrees with the NISS
severity for 30 of \LiveNissScored{} scored techniques, so it should not be read as either
a CVSS or a NISS result.

\begin{table}[H]
\centering
\caption{TARA filed severity distribution (v\SnapshotRegistryVersion{} snapshot; hand-assigned).}
\label{tab:severity}
\small
\begin{tabular}{@{}l r r@{}}
\toprule
\textbf{Severity} & \textbf{Count} & \textbf{Percentage} \\
\midrule
Critical & 29 & 26.6\% \\
High     & 60 & 55.0\% \\
Medium   & 17 & 15.6\% \\
Low      &  3 &  2.8\% \\
\bottomrule
\end{tabular}
\end{table}

The weighting toward high and critical reflects the author's judgement at filing, not a
computed result.

\subsection{Category Breakdown}

The \SnapshotTechniques{} techniques distribute across eight operational categories:

\begin{table}[H]
\centering
\caption{TARA techniques by operational category.}
\label{tab:categories}
\small
\begin{tabular}{@{}l l r@{}}
\toprule
\textbf{ID} & \textbf{Category} & \textbf{Count} \\
\midrule
SE & Signal Eavesdropping    & 20 \\
CI & Cognitive Interference  & 18 \\
EX & Data Exfiltration       & 19 \\
DM & Data Manipulation       & 15 \\
SI & Signal Injection        & 13 \\
PE & Privilege Escalation    &  9 \\
DS & Denial of Service       &  8 \\
PS & Physical Safety         &  7 \\
\bottomrule
\end{tabular}
\end{table}

\subsection{Dual-Use Mapping}

TARA assesses each technique for therapeutic analogs. The same physical mechanisms that
enable attacks---electromagnetic stimulation, signal decoding, neuromodulation---are
the mechanisms underlying established therapies such as deep brain stimulation
(DBS)~\cite{lozano2019dbs}, transcranial magnetic stimulation
(TMS)~\cite{hallett2007tms}, and neurofeedback.

\begin{table}[H]
\centering
\caption{Dual-use classification of \SnapshotTechniques{} TARA techniques.}
\label{tab:dualuse}
\small
\begin{tabular}{@{}l l r r@{}}
\toprule
\textbf{Classification} & \textbf{Definition} & \textbf{Count} & \textbf{\%} \\
\midrule
Confirmed    & Published clinical use exists      & 53 & 48.6\% \\
Probable     & Under active clinical investigation & 16 & 14.7\% \\
Possible     & Theoretical therapeutic mapping     &  9 &  8.3\% \\
Silicon Only & No tissue analog; purely digital    & 25 & 22.9\% \\
Not Mapped   & Dual-use assessment pending         &  6 &  5.5\% \\
\bottomrule
\end{tabular}
\end{table}

Of the \SnapshotTechniques{} techniques, 78 (71.6\%) have a confirmed, probable or
possible therapeutic analog, and 25 are silicon-only. The finding is that many attack
mechanisms are shared with therapies, not that every attack has a therapeutic counterpart:
the mechanism is shared, and the boundary between therapy and attack is consent, dosage and
oversight. Defences cannot simply remove a capability that treatment depends on.

\subsection{Representative Techniques}

Table~\ref{tab:representative} shows five registry entries spanning categories and dual-use
classes. In v1.5 this table listed descriptions and identifiers that did not match the
registry entries carrying those identifiers (for example, QIF-T0015 was described as a P300
side channel but is directed-energy thermal damage); it now quotes the registry.

\begin{table}[H]
\centering
\caption{Five representative TARA entries as stored in the registry (status: snapshot
$\to$ live).}
\label{tab:representative}
\small
\begin{tabularx}{\textwidth}{@{}l X l l X@{}}
\toprule
\textbf{ID} & \textbf{Name} & \textbf{Status} & \textbf{Dual-use} & \textbf{Therapeutic analog} \\
\midrule
T0001 & Signal injection & Confirmed $\to$ Theoretical & Confirmed & tDCS/tACS neuromodulation \\
T0035 & P300 interrogation & Demonstrated & Confirmed & P300 BCI for locked-in patients \\
T0058 & Calibration poisoning & Demonstrated & Probable & Adaptive BCI calibration \\
T0072 & Transducer inversion (speaker-to-microphone) & Confirmed $\to$ Demonstrated & Silicon only & --- \\
T0090 & WiFi CSI passive body sensing & Demonstrated & Confirmed & Contactless vital-sign monitoring \\
\bottomrule
\end{tabularx}
\end{table}


% ═══════════════════════════════════════════════════════════════
% Section 5: NISS — Neural Impact Scoring System
% ═══════════════════════════════════════════════════════════════

\section{NISS: Neural Impact Scoring System}
\label{sec:niss}

The Neural Impact Scoring System (NISS) is a proposed CVSS v4.0 extension
designed to conform with FIRST.org's official extension mechanism (User Guide
\S3.11)~\cite{first2023cvss4userguide}. NISS adds six metrics that capture
dimensions CVSS was never designed to express: biological tissue damage, cognitive
reconnaissance, cognitive disruption, consent boundary violations, damage
reversibility, and neuroplastic consequences.

NISS scores impact only; it has no likelihood or exploitability term and is meant to sit
beside a CVSS base score, not to replace it. It is also the technique-level input to a
proposed multi-layer Neurosecurity Score (NSv2.1), which is unvalidated and outside the
scope of this paper.

\subsection{Gap Analysis: Why CVSS Alone Is Insufficient}

Each TARA technique carries a CVSS v4.0 base vector and a gap group recording how much of
its impact CVSS fails to capture. Table~\ref{tab:gap} gives the counts published in v1.5
beside those recomputed from the registry file committed with v1.5.

\begin{table}[H]
\centering
\caption{CVSS v4.0 gap groups across the \SnapshotTechniques{} snapshot techniques: as
published in v1.5, and recomputed from the committed snapshot file.}
\label{tab:gap}
\small
\begin{tabularx}{\textwidth}{@{}c X r r@{}}
\toprule
\textbf{Group} & \textbf{Gap Description} & \textbf{v1.5 text} & \textbf{Snapshot file} \\
\midrule
1 & CVSS captures most impact; NISS adds nuance        & 13 & 14 \\
2 & CVSS captures exploitability but misses part of the impact & 31 & 30 \\
3 & CVSS cannot express the primary impact             & 64 & 59 \\
--- & No gap group assigned                            & 1  & 6 \\
\midrule
  & \textbf{Groups 1--3} & \textbf{\SnapshotGapTechniques{} (\SnapshotGapPercent\%)} & \textbf{103 (94.5\%)} \\
\bottomrule
\end{tabularx}
\end{table}

The v1.5 counts could not be reproduced from the registry file committed with that
version, which gives 59 techniques (54.1\%) in Group~3 and 103 (94.5\%) in any group; this
revision uses the reproducible counts in its text and retains the published ones in the
table for the record. Group~1 techniques are mostly captured by CVSS, so the share that
\emph{needs} NISS is better read as Groups~2 and~3, 89 techniques (81.7\%). In Group~3 the
most severe consequence is tissue damage, disruption of neural function or lasting
harm---dimensions for which CVSS has no metric. Live: \LiveGapTechniques{} of
\LiveTechniques{} (\LiveGapPercent\%) carry a gap group.

\subsection{Extension Metrics}

NISS defines six extension metrics, each with a graduated value set. The
metrics are designed to be orthogonal to CVSS base metrics: they capture impact
dimensions that exist only because the target system interfaces with biological
neural tissue.

\subsubsection{BI: Biological Impact}

Direct harm to neural tissue, organs, or physiological function. This dimension
has no equivalent in CVSS.

\begin{table}[H]
\centering
\small
\begin{tabular}{@{}l l l p{7cm}@{}}
\toprule
\textbf{Value} & \textbf{Label} & \textbf{Score} & \textbf{Description} \\
\midrule
N & None     & 0.0  & No tissue interaction or physical harm \\
L & Low      & 3.3  & Temporary discomfort, minor sensory disruption, reversible tissue stress \\
H & High     & 6.7  & Significant tissue damage, seizure induction, involuntary motor activation \\
C & Critical & 10.0 & Life-threatening or permanently disabling neural harm \\
\bottomrule
\end{tabular}
\end{table}

\subsubsection{CR: Cognitive Reconnaissance}

Degree of unauthorised read access: decoding of speech or motor intent by models trained
on a cooperating user, inference from neural data, intent extraction. CVSS treats such
reads only as loss of data confidentiality. CR maps to the Cognitive Liberty and Mental
Privacy neurorights dimensions.

\begin{table}[H]
\centering
\small
\begin{tabular}{@{}l l l p{7cm}@{}}
\toprule
\textbf{Value} & \textbf{Label} & \textbf{Score} & \textbf{Description} \\
\midrule
N & None     & 0.0  & No cognitive read access \\
L & Low      & 3.3  & Partial intent inference, coarse emotional state detection \\
H & High     & 6.7  & Decoding of intent, identity inference, or detailed neural data extraction \\
C & Critical & 10.0 & Sustained, broad inference from neural data in real time \\
\bottomrule
\end{tabular}
\end{table}

\subsubsection{CD: Cognitive/Functional Disruption}

Degree of unauthorized disruption to cognitive processing, sensory perception,
motor output, or autonomic regulation---including perception manipulation,
identity modification, cognitive coercion, and false stimuli injection. CD maps
to the Mental Integrity neurorights dimension and extends to ICD-10-CM chapters
G~(nervous system), H~(eye/ear), and R~(symptoms/signs).

\begin{table}[H]
\centering
\small
\begin{tabular}{@{}l l l p{7cm}@{}}
\toprule
\textbf{Value} & \textbf{Label} & \textbf{Score} & \textbf{Description} \\
\midrule
N & None     & 0.0  & No cognitive write access or alteration \\
L & Low      & 3.3  & Minor perceptual distortion, detectable interference \\
H & High     & 6.7  & Perception manipulation, involuntary behavioral changes \\
C & Critical & 10.0 & Cognitive coercion, identity modification, or complete loss of cognitive autonomy \\
\bottomrule
\end{tabular}
\end{table}

\subsubsection{CV: Consent Violation}

Degree of violation of informed consent or cognitive autonomy. Ordered by
severity: covert (implicit) violations are worse than detectable (explicit) ones.

\begin{table}[H]
\centering
\small
\begin{tabular}{@{}l l l p{7cm}@{}}
\toprule
\textbf{Value} & \textbf{Label} & \textbf{Score} & \textbf{Description} \\
\midrule
N & None              & 0.0  & Operating within explicitly consented boundaries \\
P & Partial           & 3.3  & Action exceeds scope but subject retains some awareness \\
E & Explicit          & 6.7  & Direct violation of consent boundaries, but detectable \\
I & Implicit (covert) & 10.0 & Covert manipulation the patient cannot detect or refuse \\
\bottomrule
\end{tabular}
\end{table}

\subsubsection{RV: Reversibility}

Whether the damage can be undone. IT assets can be restored from backup. Neural
tissue cannot be rebooted.

\begin{table}[H]
\centering
\small
\begin{tabular}{@{}l l l p{7cm}@{}}
\toprule
\textbf{Value} & \textbf{Label} & \textbf{Score} & \textbf{Description} \\
\midrule
F & Full        & 0.0  & Effects fully reverse when attack stops \\
T & Temporary   & 3.3  & Effects reverse over hours to days \\
P & Partial     & 6.7  & Some effects permanent, some reversible \\
I & Irreversible & 10.0 & Permanent neural tissue destruction or cognitive change \\
\bottomrule
\end{tabular}
\end{table}

\subsubsection{NP: Neuroplasticity}

Whether the attack exploits or induces neuroplastic changes---the brain's ability
to rewire itself. This has no digital equivalent.

\begin{table}[H]
\centering
\small
\begin{tabular}{@{}l l l p{7cm}@{}}
\toprule
\textbf{Value} & \textbf{Label} & \textbf{Score} & \textbf{Description} \\
\midrule
N & None       & 0.0  & No neuroplastic effect \\
T & Temporary  & 3.3  & Short-term synaptic changes that decay within hours--weeks \\
P & Partial    & 6.7  & Moderate synaptic reorganization persisting weeks to months \\
S & Structural & 10.0 & Long-term or permanent neural pathway changes \\
\bottomrule
\end{tabular}
\end{table}

\subsection{PINS Flag}

NISS introduces the Potential Impact to Neural Safety (PINS) flag---a
binary indicator triggered when:

\begin{equation}
\text{PINS} = \begin{cases}
  \texttt{true} & \text{if } \text{BI} \geq \text{High} \;\lor\; \text{RV} = \text{Irreversible} \\
  \texttt{false} & \text{otherwise}
\end{cases}
\end{equation}

A PINS flag calls for safety review regardless of overall score. In the snapshot, 33 of
\SnapshotTechniques{} techniques are PINS-flagged (30.3\%); in the live registry,
\LivePinsFlagged{} of the \LiveNissScored{} scored techniques are.

\subsection{Scoring Formula}

The NISS score is computed as the weighted mean of the six metric scores, rounded up to
one decimal place:

\begin{equation}
\text{NISS} = \frac{w_{\text{BI}} \cdot \text{BI} + w_{\text{CR}} \cdot \text{CR} + w_{\text{CD}} \cdot \text{CD} + w_{\text{CV}} \cdot \text{CV} + w_{\text{RV}} \cdot \text{RV} + w_{\text{NP}} \cdot \text{NP}}
              {w_{\text{BI}} + w_{\text{CR}} + w_{\text{CD}} + w_{\text{CV}} + w_{\text{RV}} + w_{\text{NP}}}
\label{eq:niss}
\end{equation}

In the default profile, the weights are $w_{\text{BI}} = w_{\text{CV}} =
w_{\text{RV}} = w_{\text{NP}} = 1.0$ and $w_{\text{CR}} = w_{\text{CD}} = 0.5$.
This normalization ensures that the split of the former Cognitive Integrity (CG)
metric into Cognitive Reconnaissance (CR) and Cognitive Disruption (CD) in v1.1
does not inflate the cognitive dimension's influence: at $0.5 + 0.5 = 1.0$ out of
a total weight sum of~5.0, cognition retains its original 20\% share from the
five-metric v1.0 formula. The CR/CD split reflects the fundamental distinction
between read attacks (inference from neural data) and write attacks
(perception manipulation, cognitive coercion)---attack vectors with different
mitigations that map to different neurorights dimensions.
NISS supports four context profiles with differential weights that override
the default normalization:

\begin{itemize}
  \item \textbf{Clinical}: BI, CD and RV at 2.0, others at 1.0
  \item \textbf{Research}: CR and CV at 2.0, CD and NP at 1.5
  \item \textbf{Consumer}: CR and CV at 2.0, others at 1.0
  \item \textbf{Military}: BI, CR and CD at 2.0, RV and NP at 1.5, CV at 0.5
\end{itemize}

\noindent All stored scores use the default weights. The military profile halves the
weight of consent, a choice with ethical content that has not been reviewed.

\subsection{Vector Format}

The NISS vector rides alongside the CVSS v4.0 base vector:

\begin{lstlisting}
CVSS:4.0/AV:N/AC:L/AT:N/PR:N/UI:N/
  VC:H/VI:H/VA:H/SC:N/SI:N/SA:N
NISS:1.1/BI:H/CR:H/CD:C/CV:I/RV:P/NP:S
\end{lstlisting}

This dual-vector architecture means security teams can triage using familiar CVSS
scores while BCI-specific teams see the neural dimensions that determine whether a
vulnerability is a software bug or a patient safety emergency.

\subsection{Scoring States and Stated Limits}
\label{sec:niss-limits}

A NISS score is six judgements about impact. Scoring a technique whose sources support only
one of them means inventing the other five, so the registry now distinguishes three
\emph{scoring states}:

\begin{itemize}
  \item \textbf{Scored}: all six metrics are supported by the technique's sources
        (\LiveNissScored{} live techniques; these were scored before the rule was
        introduced and have not yet been re-examined against it).
  \item \textbf{Provisional}: some metrics are supported; the record stores those
        values, a source for each, and the lowest and highest score the formula could
        return (\LiveNissProvisional{}). A range of 0.4 to 9.4 says that one metric is
        known and five are not; it is a bound, not an estimate, and not a severity.
  \item \textbf{Unscored}: no peer-reviewed source, or no metric supported
        (\LiveNissUnscored{}).
\end{itemize}

\noindent Read the scoring state before the score. Three limits follow from the formula and
the data:

\begin{enumerate}
  \item \textbf{Impact only.} NISS has no likelihood term, so a technique with several
        disclosed vulnerabilities can score low.
  \item \textbf{Read-only harm cannot exceed 3.0.} With CR at Critical, CV at Implicit
        and nothing physical, the score is $(0.5 \cdot 10 + 10)/5 = 3.0$. NISS therefore
        rates every pure privacy-and-consent harm as low, whatever its scale. This is a
        property of the default weights, and it is the clearest reason to report the
        vector, not only the number.
  \item \textbf{No calibration or inter-rater study.} The levels were assigned by one
        author with AI assistance (Section~\ref{sec:limitations}).
\end{enumerate}

\noindent NISS scores signal-level interference with a person's neural function and the
consent under which it occurs; it does not measure mental states, and a score is not a
clinical or legal finding. Whether NISS should be reported as a single number or as its
vector with the PINS flag is an open question.

\subsection{NISS Severity Distribution}

The NISS severity distribution in the snapshot, as published in v1.5:

\begin{table}[H]
\centering
\caption{NISS severity distribution (\SnapshotTechniques{} snapshot techniques, as published in v1.5).}
\label{tab:niss-severity}
\small
\begin{tabular}{@{}l r r@{}}
\toprule
\textbf{NISS Severity} & \textbf{Count} & \textbf{\%} \\
\midrule
High     & 19 & 17.4\% \\
Medium   & 37 & 33.9\% \\
Low      & 52 & 47.7\% \\
None     &  1 &  0.9\% \\
\bottomrule
\end{tabular}
\end{table}

The registry file committed with v1.5 gives 21 high, 37 medium, 50 low and 1 none, and
several of its stored scores do not reproduce from Equation~\ref{eq:niss}; the live scores
do. The NISS distribution sits lower than the filed severity (Table~\ref{tab:severity})
because NISS separates purely digital impact from impact on a person.


\subsection{Neurological Extension (v1.1)}

NISS v1.0 mapped technique outcomes exclusively to DSM-5-TR psychiatric
categories. This created a systematic blind spot: neurological disruptions---loss
of smell, tinnitus, vestibular dysfunction, motor tremor, autonomic
dysregulation---have no DSM-5-TR entry because the DSM is a psychiatric manual,
not a neurological reference.

NISS v1.1 extends the impact chain to 42~neurological conditions drawn from
ICD-10-CM chapters G~(diseases of the nervous system), H~(diseases of the eye
and ear), and R~(symptoms and signs not elsewhere classified), using Adams \&
Victor's \textit{Principles of Neurology} (12th~ed.)~as the primary clinical
reference. These conditions are organized into seven categories:

\begin{itemize}
  \item \textbf{Sensory} (14 conditions): auditory (tinnitus, hearing loss),
    visual (scotoma, cortical blindness), vestibular (vertigo, nystagmus),
    olfactory (anosmia, parosmia), gustatory (ageusia, dysgeusia), and
    somatosensory (paresthesia, allodynia)
  \item \textbf{Motor} (8 conditions): tremor, dystonia, ataxia, spasticity,
    fasciculations, myoclonus, bradykinesia, dyskinesia
  \item \textbf{Autonomic} (6 conditions): orthostatic hypotension,
    hyperhidrosis, cardiac arrhythmia, thermoregulatory dysfunction,
    gastroparesis, neurogenic bladder
  \item \textbf{Consciousness} (4 conditions): syncope, absence seizures,
    altered sensorium, coma
  \item \textbf{Speech/Language} (4 conditions): expressive aphasia, receptive
    aphasia, dysarthria, alexia
  \item \textbf{Reflex} (3 conditions): hyperreflexia, areflexia, abnormal
    plantar response
  \item \textbf{Pain} (3 conditions): neuropathic pain, trigeminal neuralgia,
    central pain syndrome
\end{itemize}

Each condition maps to specific QIF hourglass bands, neural pathway
dependencies, and NISS metrics. The CD metric was broadened from ``Cognitive
Disruption'' to ``Cognitive/Functional Disruption'' to encompass sensory, motor,
and autonomic disruption alongside cognitive interference. This definitional
change preserves backward compatibility: all existing CD scores remain valid,
and the scoring formula is unchanged. The extension adds coverage for
neurological outcomes that were previously invisible to the scoring system.

The neurological conditions are queryable via KQL as the
\texttt{neurological\_conditions} table, enabling cross-correlation with TARA
techniques, neural pathways, and hourglass bands.

% ═══════════════════════════════════════════════════════════════
% Section 6: Neural Impact Chain
% ═══════════════════════════════════════════════════════════════

\section{Neural Impact Chain}
\label{sec:nic}

The Neural Impact Chain (NIC) is a six-stage threat-modeling method that associates
attack techniques with DSM-5-TR and ICD-10-CM diagnostic category
references~\cite{apa2022dsm5tr}. The codes are \textbf{diagnostic category references for
threat modeling, not diagnostic claims}: an association says that a technique and a
condition category share an address and a function, not that the technique causes the
condition. The question the NIC helps a threat modeler ask is \emph{which clinical
categories should a review of this technique consider?} The mappings have not been reviewed
by clinicians (Section~\ref{sec:limitations}).

\subsection{Pipeline Architecture}

The NIC traces each technique through six stages:

\begin{enumerate}
  \item \textbf{Technique}: The TARA attack technique (e.g., cortical signal
        injection)
  \item \textbf{Address}: The code(s) the technique is tagged at (e.g., N7
        Neocortex, N6 Allocortex; version 1 memberships in the snapshot)
  \item \textbf{Neural Structure}: The anatomical structure at risk (e.g.,
        hippocampus, prefrontal cortex, amygdala)
  \item \textbf{Cognitive Function}: The function that structure supports (e.g.,
        memory consolidation, executive control, emotional regulation)
  \item \textbf{NISS Score}: The neural impact score, particularly the BI, CR,
        CD, and NP metrics that correlate with clinical outcomes
  \item \textbf{Category reference}: The DSM-5-TR/ICD-10-CM code(s) for the
        condition category most closely associated with disruption of that function
\end{enumerate}


\subsection{NISS-to-DSM Bridge}

A rule assigns each technique to a DSM-5-TR category cluster according to which NISS
metric dominates its profile. The rule is a modeling choice, not a clinical finding:

\begin{itemize}
  \item \textbf{BI-driven} $\to$ Motor/Neurocognitive cluster: movement and
        neurocognitive categories
  \item \textbf{CR/CD-driven} $\to$ Cognitive/Psychotic cluster: perceptual and
        cognitive categories
  \item \textbf{CV-driven} $\to$ Mood/Trauma cluster: trauma- and stressor-related
        categories
  \item \textbf{NP/RV-driven} $\to$ Persistent/Personality cluster: categories of
        lasting change
  \item \textbf{No neural impact} $\to$ Non-diagnostic: Silicon-only
        techniques with no category association
\end{itemize}

\subsection{Coverage Statistics}

All \SnapshotTechniques{} TARA techniques have been mapped through the NIC pipeline:

\begin{table}[H]
\centering
\caption{Neural Impact Chain mapping results across \SnapshotTechniques{} techniques.}
\label{tab:nic-stats}
\small
\begin{tabular}{@{}l r@{}}
\toprule
\textbf{Metric} & \textbf{Value} \\
\midrule
Techniques mapped                 & 103 / \SnapshotTechniques{} (94.5\%) \\
Unique category codes             & 29 \\
Diagnostic clusters               & 5 \\
Direct association                 & 52 (47.7\%) \\
Indirect association               & 9 (8.3\%) \\
No association (silicon-only)      & 42 (38.5\%) \\
\bottomrule
\end{tabular}
\end{table}

\subsection{Diagnostic Cluster Distribution}

The five diagnostic clusters and their technique counts:

\begin{table}[H]
\centering
\caption{DSM-5-TR diagnostic category clusters (snapshot), for threat modeling.}
\label{tab:clusters}
\small
\begin{tabularx}{\textwidth}{@{}l r X@{}}
\toprule
\textbf{Cluster} & \textbf{Count} & \textbf{Representative DSM-5-TR Codes} \\
\midrule
Non-diagnostic          & 42 & --- (silicon-only techniques) \\
Mood/Trauma             & 21 & F43.10 (PTSD), F32.9 (MDD), F44.9 \\
Cognitive/Psychotic     & 17 & F06.0 (psychosis), R41.3 (cognitive decline) \\
Motor/Neurocognitive    & 16 & G25.9 (movement disorder), G31.84 \\
Persistent/Personality  &  7 & F07.0 (personality change) \\
\bottomrule
\end{tabularx}
\end{table}

The Mood/Trauma cluster is the largest diagnostic category cluster (21 techniques),
because the bridge rule sends consent-dominated techniques there. That is a consequence of
the rule, not an observation about outcomes in people.

\subsection{Association Classes}

Each technique is placed in one of three association classes:

\begin{itemize}
  \item \textbf{Direct} (52 techniques, 47.7\%): the technique acts on the
        structure or function associated with the category. Example: unauthorised
        cortical stimulation is associated with seizure-related categories.
  \item \textbf{Indirect} (9 techniques, 8.3\%): the association runs through a
        downstream effect. Example: exfiltration of neural data is associated with
        anxiety-related categories through the effect of the disclosure, not through
        the neural signal.
  \item \textbf{None} (42 techniques, 38.5\%): Silicon-only techniques with no
        direct neural interaction and thus no category association.
\end{itemize}

\subsection{Example Walkthrough}

Consider \textbf{QIF-T0001, Signal injection}, as stored in the live registry. (The v1.5
walkthrough gave a vector, a score of 8.7 and three codes that did not match the registry
entry; it is replaced here.)

\begin{enumerate}
  \item \textbf{Technique}: QIF-T0001 (signal injection at the electrode-tissue interface)
  \item \textbf{Address}: I0, N1
  \item \textbf{Mechanism}: electrical current delivered at the interface, modulating
        local field potentials
  \item \textbf{NISS}: \texttt{BI:H/CR:H/CD:H/CV:E/RV:P/NP:T} $\to$ 6.1 (Medium), PINS
        set because BI is High
  \item \textbf{Category references}: F43.2, F45, F44.4 (primary); F82 (secondary)
  \item \textbf{Association class}: direct
\end{enumerate}

\noindent Its evidence status is theoretical: the technique carried \emph{confirmed} until
2026-10-16, when its sources were found not to demonstrate signal injection
(Section~\ref{sec:tara}).

\subsection{Impact Chains as Associations (Live Dataset)}

The live impact-chains dataset (5,546 rows) is built by joins: each technique to every
structure sharing one of its addresses, then to every pathway touching that structure, then
to that pathway's condition codes. A row therefore records a shared address, not that the
technique reaches the structure. The join currently yields rows for 91 of the
\LiveTechniques{} techniques and 10 structures. A technique with no row has no association
recorded, which is not evidence that it has no effect. Because the joins run on addresses,
the move to addressing version 2 (Section~\ref{sec:hourglass}) changes them; the snapshot
figures above use version 1.


% ═══════════════════════════════════════════════════════════════
% Section 7: Governance & Neuroethics
% ═══════════════════════════════════════════════════════════════

\section{Governance and Neuroethics}
\label{sec:governance}

This section sets out why a security framework for neural devices needs neuroethics, where
neuroethics enters the scoring and the labels, and how the framework relates to policy
instruments. None of it has been examined by an ethicist, clinician or review board.

\subsection{Why Neuroethics}
\label{sec:why-neuroethics}

Neuroethics asks what may rightly be done to and with a human brain. Roskies divided the
field into the ethics of neuroscience, which asks how brain research and technology should
be conducted and used, and the neuroscience of ethics~\cite{roskies2002neuroethics}; this
framework is concerned only with the first.

A security framework needs it because the harms of a compromised neural device are not all
harms to systems or data. Consent to a medical device can be partial and can be withdrawn;
some effects on tissue or neural function cannot be undone; and a recording of brain
activity cannot be reissued. Conventional scoring has no vocabulary for these, and
neuroethics supplies one. Ienca and Andorno~\cite{ienca2017neurorights} proposed four
rights---cognitive liberty, mental privacy, mental integrity and psychological
continuity---and Yuste et al.~\cite{yuste2017four} set out ethical priorities for
neurotechnology and artificial intelligence.

The rights are contested. Bublitz argues against inflating the catalogue of rights, while
describing Ienca and Andorno's four as better grounded in existing
scholarship~\cite{bublitz2022freedom}; Ienca's own later review warns against the same
inflation~\cite{ienca2021neurorights}; and mental privacy and mental integrity lack agreed
operational definitions~\cite{kellmeyer2022neurorights}. QIF therefore uses the four rights
as \emph{labels} for what a technique puts at stake, defines only measurable signal
properties, and does not argue that new law is needed.

Neuroethics also limits what a security framework may claim. The project holds eight
guardrails as data (\texttt{datalake/qif-guardrails.json}), each pairing a caution from the
literature with a limit on a component. Among them: activity in a brain region does not
identify a cognitive process~\cite{poldrack2006neuroimaging}, so TARA catalogues physical
interference, not cognitive content; brain data is given more credibility than it has
earned~\cite{racine2005neuroethics}; present devices do not read
thoughts~\cite{ienca2018brainleaks,wexler2019neurohype}, so every technique separates what
is possible now from what is projected; and neuroscience is dual-use, pursued by national
security agencies~\cite{tennison2012neurorights}, so every threat description is defensive
and paired with a control. The evidence rule, the three scoring states and the labelling of
diagnostic codes as category references (Sections~\ref{sec:tara}--\ref{sec:nic}) apply
these guardrails in practice.

\subsection{How Neuroethics Enters the Score}

Three commitments are fixed in the NISS arithmetic (Equation~\ref{eq:niss}):

\begin{enumerate}
  \item \textbf{Consent is a scored metric at full weight.} A technique that changes
        nothing physically still scores if it overrides consent, although it cannot
        score above 3.0 (Section~\ref{sec:niss-limits}). Of the \LiveNissScored{} live
        scored techniques, 68 are rated Explicit, 49 Implicit, 5 Partial and 7 None on CV.
  \item \textbf{Harm that cannot be undone counts for more.} Reversibility and
        Neuroplasticity carry full weight, expressing non-maleficence as a number, and the
        PINS flag applies to any irreversible technique whatever its total.
  \item \textbf{Reading and disruption are scored separately.} CR and CD carry half weight
        each, so together they weigh as much as one other metric.
\end{enumerate}

\noindent Rights labels are then assigned by rule: High or Critical CR earns mental privacy
and cognitive liberty; High or Critical CD or BI earns mental integrity; further labels come
from a technique's category, its category references and short hand-made lists. In the live
registry 142 of \LiveTechniques{} techniques carry at least one label (mental privacy 111,
mental integrity 102, cognitive liberty 87, psychological continuity 79). The labels do not
feed back into the score.

\subsection{Consent Tiers}

Each technique's governance record names a consent tier:

\begin{table}[H]
\centering
\caption{Consent tiers: snapshot (v\SnapshotRegistryVersion{}) and live counts.}
\label{tab:consent}
\small
\begin{tabularx}{\textwidth}{@{}l X r r@{}}
\toprule
\textbf{Tier} & \textbf{Requirement} & \textbf{v\SnapshotRegistryVersion{}} & \textbf{Live} \\
\midrule
Standard   & Standard clinical consent                         & 22 & 42 \\
Enhanced   & Extended disclosure and monitoring                & 51 & 70 \\
IRB        & Full institutional review board / ethics committee review & 28 & 52 \\
Prohibited & Prohibited outside controlled research            & 2  & 9 \\
---        & Not assigned                                      & 6  & 0 \\
\bottomrule
\end{tabularx}
\end{table}

\subsection{What an Institutional Programme Has, and QIF Lacks}

The guardrails above are self-imposed: no one outside the project checks them. The
neuroethics programme of the NIH BRAIN Initiative is a published example of more. It treats
neuroethics as concurrent with tool development rather than as review after the
fact~\cite{ramos2017brainneuroethics}, and rests on four standing activities: a Neuroethics
Working Group of neuroethics and neuroscience experts that advises the programme; funded
neuroethics research; scanning of the research portfolio for the ethical questions it
raises; and workshops on topics that steer the programme~\cite{bianchi2018brainneuroethics}.
A later account describes this as integrating neuroethics into neuroscience rather than
appending it~\cite{ramos2019brainintegrating}.

QIF has none of the first, second or fourth, and performs the third only as the author
applying the guardrails to the author's own work. This is a gap statement, not a plan; it
keeps the four activities visible as the standard the framework falls short of. Nothing here
implies that the NIH, the BRAIN Initiative or its working group has reviewed, endorsed or has
any connection to QIF.

\subsection{Alignment with Policy Instruments}

UNESCO's Recommendation on the Ethics of Neurotechnology~\cite{unesco2025recommendation},
adopted in November 2025, sets out values, principles and policy action areas.
Table~\ref{tab:unesco} maps its elements to QIF components. The mapping is the author's
self-assessment; it has not been reviewed by UNESCO or any third party, and ``mapped'' means
that a component addresses the element, not that the element is met.

\begin{table}[H]
\centering
\caption{Self-assessed mapping of UNESCO Recommendation elements to QIF components.}
\label{tab:unesco}
\small
\begin{tabularx}{\textwidth}{@{}l l l X@{}}
\toprule
\textbf{UNESCO Element} & \textbf{Type} & \textbf{Status} & \textbf{QIF Component} \\
\midrule
Human rights \& dignity        & Value     & Mapped  & Rights labels, consent tiers \\
Health \& well-being           & Value     & Mapped  & PINS flag, biological impact scoring \\
Diversity                      & Value     & Mapped  & Open-source, multi-stakeholder model \\
Sustainability                 & Value     & Mapped  & Apache 2.0 license, post-trial access \\
Professional integrity         & Value     & Mapped  & Transparency audit trail \\
Proportionality                & Principle & Mapped  & Graduated severity scoring \\
Freedom of thought             & Principle & Mapped  & Cognitive liberty label, consent metric \\
Privacy                        & Principle & Mapped  & Cognitive reconnaissance (CR) / disruption (CD) metrics \\
Protection of children         & Principle & Mapped  & Age-tiered assent framework (governance documents) \\
Consumer protection            & Policy    & Mapped  & Default-deny architecture (governance documents) \\
Enhancement regulation         & Policy    & Partial & Technical infrastructure only \\
Workplace protections          & Policy    & Mapped  & Mental privacy scoring \\
Behavioral influence           & Policy    & Mapped  & Consent violation metric (CV) \\
Health and well-being          & Policy    & Mapped  & NISS scoring, PINS flag \\
Oversight \& governance        & Policy    & Mapped  & Regulatory mapping (below) \\
Access \& equity               & Policy    & Mapped  & Open-source release \\
\bottomrule
\end{tabularx}
\end{table}

\noindent The table lists 16 rows. The v1.5 text described them as 15 of 17 elements
``implemented'' with two partial, which did not match its own table; ``implemented'' is
replaced by ``mapped'' because no row has been tested against the Recommendation. Two
components cited (the age-tiered assent framework and default-deny architecture) are
described in the project's governance documents and are not part of the TARA or NISS data
analysed here.

\subsection{Regulatory Mapping}

The framework relates to existing instruments as follows. None requires TARA or NISS, and
no regulator has reviewed either.

\begin{itemize}
  \item \textbf{FDA Section 524B}~\cite{fda524b}: mandatory cybersecurity documentation
        for connected medical devices, including threat modeling; TARA could inform the
        threat model and NISS could sit beside CVSS in it
  \item \textbf{EU MDR 2017/745}~\cite{eumdr2017}: medical device risk management, for
        which TARA offers a neural threat taxonomy
  \item \textbf{ISO 14971}~\cite{iso14971}: risk management for medical devices; the NIC
        offers a structured association from technical failure to categories of harm
  \item \textbf{HIPAA}~\cite{hipaa1996} and \textbf{GDPR}~\cite{gdpr2016}: health and
        sensitive personal data, under which neural data may fall (GDPR Article 9)
  \item \textbf{Colorado Privacy Act}~\cite{colorado2024neuraldata}: amended in 2024 to
        classify neural data as sensitive data
\end{itemize}

% ═══════════════════════════════════════════════════════════════
% Section 8: Case Studies
% ═══════════════════════════════════════════════════════════════

\section{Case Studies}
\label{sec:case-studies}

To show the practical difference between CVSS-only and CVSS+NISS scoring, we score four
registry entries with both systems, one from each gap group plus a read-only case.

\textbf{Correction.} In v1.5, Cases 1--4 used identifiers whose registry entries describe
different techniques (QIF-T0015, T0034 and T0050 are directed-energy damage, working-memory
poisoning and hardware fault injection, not the P300, calibration and surveillance attacks
described), and their NISS vectors and scores did not match the registry. Some case scores
were also computed as an unweighted mean of six, which is not Equation~\ref{eq:niss}. This
revision replaces them with entries as stored in the live registry.

\textbf{Method and caveat.} NISS scores are recomputed from the stored vectors with
Equation~\ref{eq:niss} and match the stored values. CVSS v4.0 base scores were computed
from the stored base vectors with the open-source \texttt{cvss} Python library (v3.6) and
have not been checked against the FIRST.org calculator. The vectors were unchanged between
the snapshot and the live registry, but the snapshot file stored NISS scores that do not
reproduce from the equation (for example 3.4 for QIF-T0035 against 2.7). None of these
techniques has been executed against an implanted device; Case~5 is the only analysis of a
real vulnerability.

\begin{table}[H]
\centering
\caption{CVSS v4.0 vs.\ NISS for four registry entries (live registry).}
\label{tab:case-studies}
\small
\begin{tabularx}{\textwidth}{@{}l X l c c@{}}
\toprule
\textbf{Technique} & \textbf{Name} & \textbf{Status} & \textbf{CVSS} & \textbf{NISS} \\
\midrule
QIF-T0001 & Signal injection       & Theoretical  & 5.8 (Med)  & 6.1 (Med), PINS \\
QIF-T0035 & P300 interrogation     & Demonstrated & 6.9 (Med)  & 2.7 (Low) \\
QIF-T0058 & Calibration poisoning  & Demonstrated & 7.1 (High) & 6.0 (Med) \\
QIF-T0072 & Transducer inversion   & Demonstrated & 5.6 (Med)  & 2.0 (Low) \\
\bottomrule
\end{tabularx}
\end{table}

\subsection{Case 1: Signal Injection (QIF-T0001, Gap Group 3)}

CVSS rates this Medium (5.8): the vector requires physical access and attack preconditions,
with high integrity and availability impact on the system and high integrity impact on
subsequent systems. What it cannot express is the consequence for the person: current
delivered at the electrode-tissue boundary can cause tissue stress or involuntary activation.
\texttt{NISS:1.1/BI:H/CR:H/CD:H/CV:E/RV:P/NP:T} gives 6.1 and sets PINS because BI is High.
The technique is theoretical: its two listed sources are a position paper and a P300 probing
study, so the score describes a modeled impact, not an observed one.

\subsection{Case 2: P300 Interrogation (QIF-T0035, Read-Only)}

CVSS rates this Medium (6.9) for high confidentiality impact on the device and on subsequent
systems. The leaked data is the person's own evoked responses, which reveal recognition of
stimuli such as PINs, faces or places~\cite{martinovic2012feasibility}.
\texttt{NISS:1.1/BI:N/CR:H/CD:H/CV:E/RV:F/NP:N} gives
$(0 + 0.5 \cdot 6.7 + 0.5 \cdot 6.7 + 6.7 + 0 + 0)/5 = 2.68 \to 2.7$. The low score is
a direct consequence of the read-only ceiling (Section~\ref{sec:niss-limits}); the CR and CV
values, not the total, carry the information. The registry's CD:H for a read attack is a
scoring choice that warrants review.

\subsection{Case 3: Calibration Poisoning (QIF-T0058, Gap Group 2)}

CVSS rates this High (7.1) for integrity impact during calibration. Corrupted calibration
data makes the decoder learn wrong mappings between neural signals and intended actions,
and the user may adapt to the corrupted mapping. \texttt{NISS:1.1/BI:L/CR:H/CD:H/CV:I/RV:P/NP:T}
gives $(3.3 + 3.35 + 3.35 + 10 + 6.7 + 3.3)/5 = 6.0$; CV:I records that the user cannot
detect the corruption. The entry's source is Schroder et al.~\cite{schroder2025cyberrisks}; the
registry pairs it with signed calibration data and recalibration audits as controls.

\subsection{Case 4: Transducer Inversion (QIF-T0072, Gap Group 1)}

A speaker is re-tasked as a microphone. CVSS rates this Medium (5.6), and its
confidentiality metrics describe the impact well. NISS adds little:
\texttt{NISS:1.1/BI:N/CR:N/CD:N/CV:I/RV:F/NP:N} gives 2.0, and the dual-use class is
silicon-only. The case is included to show that NISS is not meant to raise every score.

\subsection{Case 5: Real-World Vulnerability Disclosure}

The \qif framework has been applied to real vulnerability research. During
systematic analysis of the BCI software ecosystem, we identified a multi-phase
exploit chain in an open-source library used in clinical and research BCI
pipelines.

The exploit chain demonstrates escalation from device-side vulnerabilities
(S2/S3) to potential impact on clinical decision-making through corrupted data.
CVSS scores the software vulnerabilities accurately; NISS records the downstream
risk to patients whose clinical care depends on the integrity of the data stream.

Responsible disclosure is in progress. Specific vulnerability details, including
affected software and CWE identifiers, will be published after coordinated
disclosure concludes.

% ═══════════════════════════════════════════════════════════════
% Section 9: Limitations and Future Work
% ═══════════════════════════════════════════════════════════════

\section{Limitations and Future Work}
\label{sec:limitations}

We present these limitations transparently to guide future validation efforts and
to prevent overstatement of the framework's current maturity. QIF, TARA and NISS are
proposed and unvalidated: not independently peer-reviewed as a whole, and not adopted by
any standards body or device maker.

\subsection{No Empirical Validation on Real BCI Devices}

The \qif framework has not been validated against operational BCI hardware.
The TARA taxonomy was developed through literature review, threat modeling, and
systematic analysis rather than penetration testing of actual neural devices.
While the framework has been applied to one real software vulnerability
(Section~\ref{sec:case-studies}), this covers only device-side code. Validation
against interface and neural threats requires access to implanted BCI
patients and clinical environments---resources unavailable to independent
researchers.

\subsection{Reproducibility of the v\SnapshotRegistryVersion{} Figures}

For this revision the snapshot figures were recomputed from the registry file committed
with v1.5. The evidence-status, filed-severity, category, dual-use, consent-tier, PINS and
Neural Impact Chain counts reproduce exactly. Three do not: the gap groups (published
13/31/64, totalling \SnapshotGapTechniques{} or \SnapshotGapPercent\%; file 14/30/59,
totalling 103 or 94.5\%), the NISS severity distribution (published 19/37/52/1 high to none;
file 21/37/50/1), and several stored NISS scores, which do not reproduce from
Equation~\ref{eq:niss}. The cause was not identified. The text of this revision uses the
reproducible gap counts; the published ones are kept in the tables for the record. The case
studies and representative techniques of v1.5 did not match the registry at all and were
replaced (Sections~\ref{sec:tara} and~\ref{sec:case-studies}).

\subsection{Diagnostic Category Associations Not Clinically Validated}

The Neural Impact Chain associates techniques with diagnostic category references through
neuroanatomical addresses and functional neuroscience. These associations have not been
reviewed by psychiatrists, neurologists or clinical neuroscientists. They are references for
threat modeling, not diagnostic claims, and must not inform clinical decision-making.

\subsection{NISS Weights Not Calibrated}

The default profile weights BI, CV, RV and NP at 1.0 and CR and CD at 0.5
(Equation~\ref{eq:niss}). Neither these nor the four context profiles have been calibrated
against empirical data, expert elicitation or clinical outcomes, and the read-only ceiling
of 3.0 (Section~\ref{sec:niss-limits}) is a direct consequence of the default choice.
Calibration requires expert panel scoring of representative scenarios, sensitivity analysis
across weight configurations, and correlation with observed outcomes when available.

\subsection{No Interrater Reliability Study}

NISS scores in the TARA registry were assigned by a single analyst (the author) with AI
assistance. No interrater reliability study has been conducted. CVSS interrater reliability
is a known challenge~\cite{first2023cvss4}; NISS, with its novel neural-specific metrics,
likely faces greater variability. A formal study with domain experts from both
cybersecurity and neuroscience is needed.

\subsection{Evidence and Completeness}

The analysed snapshot contains \SnapshotTechniques{} techniques (v\SnapshotRegistryVersion{}),
of which 31 are theoretical and 1 speculative. The live registry holds \LiveTechniques{}
techniques and \LiveVariants{} variant across \LiveTactics{} tactics
(v\LiveRegistryVersion{}), of which 95 are theoretical. \LiveWithCvss{} carry CVSS v4.0 base
vectors; of the NISS scores \LiveNissScored{} are complete, \LiveNissProvisional{} are
provisional ranges and \LiveNissUnscored{} are unscored under the evidence rule. The
\LiveNissScored{} scored techniques were scored before that rule and have not been
re-examined against it, so the scored count is expected to fall. Live figures are not
substituted into the analysis, because every proportion in
Sections~\ref{sec:tara}--\ref{sec:governance} has a numerator computed over the snapshot;
re-pointing the denominator would produce figures that describe no actual registry state.
Current values are derivable from \texttt{datalake/qtara-registrar.json} in the project
repository. The addressing model has not been reviewed by a human neuroanatomist, and the
catalogue records techniques one person at a time, so attacks spread across a population are
outside its model (Section~\ref{sec:related-work}).

\subsection{Single-Author Bias}

The framework was developed by a single independent researcher. While multi-model AI
review was used throughout development (Claude, Gemini, ChatGPT), the architectural
decisions, scoring assignments, and clinical associations reflect a single perspective, and
AI review does not substitute for independent human review.

\subsection{AI Tool Disclosure}

In accordance with arXiv, ACM, and IEEE policies on AI-assisted research, we
disclose the following. Large language models (Claude~3.5/4, Gemini~2.0, ChatGPT-4o)
were used during the development of this framework and versions up to v1.5. Their roles, by
section:

\begin{itemize}
  \item \textbf{Literature review} (Sections~2, 3, 7): AI assisted in identifying
        relevant publications and summarizing prior work. All citations were
        independently verified by the author via DOI resolution and publisher page
        confirmation (see version history below).
  \item \textbf{Code generation} (figures, TARA data pipeline): AI generated
        Python scripts for data analysis, visualization, and figure generation. All
        code was reviewed and modified by the author.
  \item \textbf{Writing assistance} (all sections): AI provided draft text that was
        substantially rewritten, restructured, and verified by the author. For versions
        up to v1.5, the proportion of AI-generated text retained verbatim is estimated at
        less than 15\%.
  \item \textbf{Cross-validation}: AI models were used to cross-check factual claims
        across multiple systems (Claude, Gemini, ChatGPT) to reduce single-model
        bias. Disagreements between models were resolved by the author through
        primary source verification.
\end{itemize}

\noindent\textbf{This revision (v1.6).} The v1.6 changes were drafted on 2026-10-17 by Claude
Opus~5.5 (Anthropic), operating through Claude Code, from the project's data files, the v9
whitepaper and the repository history. The AI system wrote the new and rewritten text (the
abstract; Section~\ref{sec:why-neurosecurity}; the cognitive-neurosecurity subsection of
Section~\ref{sec:related-work}; Section~\ref{sec:hourglass}; the evidence, hazards and
representative-technique subsections of Section~\ref{sec:tara}; the gap analysis and
Section~\ref{sec:niss-limits}; the walkthrough and association subsections of
Section~\ref{sec:nic}; Sections~\ref{sec:why-neuroethics} and the remainder of
Section~\ref{sec:governance}; Section~\ref{sec:case-studies}; this section's reproducibility,
evidence and version notes; and Section~\ref{sec:conclusion}). It recomputed every figure
it added from the data files with scripts it ran, computed the CVSS scores in
Section~\ref{sec:case-studies}, and re-resolved each added citation through the Crossref or
arXiv API on the same date. Text new in v1.6 is predominantly AI-drafted and was pending
the author's review when drafted; it should not be read as author-reviewed until the
Zenodo version description says so.

\noindent\textbf{Human-originated contributions:} The original band model, the TARA
taxonomy structure, the NISS scoring methodology, the Neural Impact Chain method, all
diagnostic category associations, all NISS metric assignments, and the research conclusions
of v1.0--v1.5 were conceived, designed and decided by the author. The addressing revision,
the evidence rule and the retirement of the confirmed status were made in the project's data
before this revision, which reports them.

\noindent The author takes full responsibility for all content in this paper,
irrespective of how it was generated. A complete, auditable transparency log
documenting AI contributions and verification steps is maintained at
\url{https://github.com/qinnovates/qinnovate/blob/main/governance/TRANSPARENCY.md}.

\paragraph{Version history.}
An earlier version of this preprint (v1.0) contained citation errors
introduced during AI-assisted bibliography construction, including three
fabricated entries. These were corrected in v1.1 through a two-pass
independent verification audit. Version 1.2 corrected an internal percentage
inconsistency and added the author responsibility statement above.
Version~1.3 expanded the regulatory context with FDORA/PATCH Act Section~524B analysis and
added Schroder et al.~(2025) to the related work. Version~1.4 added the Sherman \&
Guillery~(2006) citation, replaced a mislabeled DBS citation with Hallett~(2007) for TMS,
corrected case study scores, fixed the dual-use classification label, and completed a
citation and claims verification pass. Version~1.5 expanded the registry from 102 to
\SnapshotTechniques{} techniques, split the NISS Cognitive Integrity metric into Cognitive
Reconnaissance and Cognitive Disruption (NISS~v1.1), normalized default weights, and added
neurorights mapping for 103 of \SnapshotTechniques{} techniques.
Version~1.6 (October 2026) replaces the hourglass stack with addressing version 2; adds the
evidence rule, the three scoring states, the retirement of the confirmed status, the hazards
dataset and NISS's stated limits; replaces case studies and representative techniques that
did not match the registry; reports which snapshot figures do and do not reproduce; adds
arguments for why neurosecurity and neuroethics are needed, with twenty newly cited
sources; qualifies or removes priority claims, universal negatives and diagnostic language;
removes six figures whose hard-coded values or labels no longer matched the tables or the
model; and condenses the future-work list.

\noindent\textbf{Note on AI use and preprint status:} I did not find anything
in Zenodo's terms of use that prohibits using AI to help write a preprint, so
I used it, heavily, to get this initial version out the door. The AI helped
me move faster, but the ideas, the architecture, the scoring, the clinical
mappings, those are mine. This preprint is me planting a seed. The real work
is what comes next: writing it properly with human collaborators, getting the
DSM-5-TR mappings reviewed by actual clinicians, running empirical validation,
and revising it until it holds up to scrutiny. There is a lot more to do.

\subsection{Future Work}

\begin{enumerate}
  \item \textbf{Re-examination of existing scores}: apply the evidence rule to the
        \LiveNissScored{} scored techniques
  \item \textbf{Expert review}: a neuroanatomist's review of the addressing model before
        the routes are integrated, and clinical review of the category associations
  \item \textbf{Interrater reliability and weight calibration}: a formal study with
        cybersecurity and neuroscience experts, with sensitivity analysis of the weights
  \item \textbf{Checks in the build}: run the band check, a score recomputation and the
        route gate in continuous integration
  \item \textbf{Standards input}: offer NISS to FIRST.org's CVSS community only after at
        least one independent researcher has reviewed the scoring method; no submission has
        been made
  \item \textbf{Empirical testing}: testing against BCI hardware in controlled environments
        with ethics approval
  \item \textbf{Exploratory directions} documented in the repository and not claims of this
        paper, each requiring ethics review before any implementation: hash-chained audit
        logs for interface firmware; an EEG dataset pipeline for validating components
        (\url{https://github.com/qinnovates/qinnovate/blob/main/shared/EEG-DATA-PIPELINE.md});
        a commodity-hardware vision-prosthesis security prototype
        (\url{https://github.com/qinnovates/qinnovate/blob/main/docs/research/lidar-bci-bridge-technical-spec.md});
        and consent-gated public-safety concepts
\end{enumerate}

% ═══════════════════════════════════════════════════════════════
% Section 10: Conclusion
% ═══════════════════════════════════════════════════════════════

\section{Conclusion}
\label{sec:conclusion}

Brain-computer interfaces are moving from laboratory prototypes to commercial medical
devices. Security frameworks designed for information technology are necessary for these
devices but not sufficient: a compromised neural device can act on a person, can expose
data that cannot be reissued, and can outlast the party that supports it. That is the case
for neurosecurity. Because the harms include violated consent and effects that cannot be
undone, and because claims about brains are easily overstated, a neurosecurity framework
also needs neuroethics, both to decide what counts as harm and to limit what the framework
itself may claim.

This paper presented the \qif framework, a proposed and unvalidated system comprising an
addressing model for what a device touches, a threat catalogue (TARA), an impact-only CVSS
v4.0 extension (NISS, six metrics, with three scoring states), and the Neural Impact Chain,
a threat-modeling method that associates techniques with diagnostic category references.
In the \SnapshotTechniques{}-technique snapshot analysed here, CVSS cannot express the
primary impact of 59 techniques and misses part of it for 30 more; 71.6\% of techniques have
a recorded therapeutic analog, showing how often the mechanism of an attack is the mechanism
of a therapy; and 56.0\% are in a direct or indirect association class.

This revision also corrected the paper. It withdrew the hourglass stack and its chokepoint,
retired an evidence status that conflated unrelated kinds of evidence, replaced case studies
that did not match the registry, and reported which published figures could not be
reproduced. The framework's limits remain: no empirical validation on BCI hardware, no
clinical review of the category associations, no expert review of the addressing model, and
single-author scoring. These are invitations to collaborate, not reasons to delay
publication.

The framework, catalogue, NISS specification and scoring data are released under the
Apache 2.0 license. We invite review and correction from the neuroscience, neuroethics,
cybersecurity and clinical communities; a finding that an entry is wrong is as useful as a
finding that it holds.


\bibliographystyle{plainnat}
\bibliography{references}

\end{document}
